% PAGES: 314-316
% Continues Theorem 10.5, opened in sec10_c1_4.tex (part (i) statement only,
% eq (10.47)-(10.48)). Supplies part (ii) of the statement below, then the
% proof, then "\end{theorem}" -- no "\begin{theorem}" or "Theorem 10.5" label
% is reprinted here.

(ii) Let $A$ be a square matrix. The gradient and Hessian of $x^tAx$ are
\begin{align}
D(x^tAx) &= \left((A+A^t)x\right)^t; \\
D^2(x^tAx) &= A+A^t.
\end{align}
\end{theorem}

\begin{proof}
(i) Let $A$ be an $n \times n$ symmetric matrix, and let $f : \R^n \to \R$ given by
\[
f(x) \;=\; x^tAx.
\]

Let $i$ fixed, with $1 \le i \le n$. From (10.20), we find that
\begin{align}
f(x) \;&=\; \sum_{1\le l,p\le n} A(l,p)x_lx_p \notag \\
&=\; A(i,i)x_i^2 \;+\; \sum_{p=1:n,p\ne i} A(i,p)x_ix_p \;+\; \sum_{l=1:n,l\ne i} A(l,i)x_lx_i \\
&\quad +\; \sum_{1\le l,p\le n,l\ne i,p\ne i} A(l,p)x_lx_p,
\end{align}

By differentiating (10.51--10.52) with respect to $x_i$, and noting that $A(i,p) =
A(p,i)$ since $A$ is a symmetric matrix, we find that
\begin{align}
\pd{f}{x_i}(x) \;&=\; 2A(i,i)x_i \;+\; \sum_{p=1:n,p\ne i} A(i,p)x_p \;+\; \sum_{l=1:n,l\ne i} A(l,i)x_l \notag \\
&=\; 2A(i,i)x_i \;+\; \sum_{p=1:n,p\ne i} A(i,p)x_p \;+\; \sum_{p=1:n,p\ne i} A(p,i)x_p \notag \\
&=\; 2A(i,i)x_i \;+\; \sum_{p=1:n,p\ne i} (A(i,p)+A(p,i))x_p \notag \\
&=\; 2A(i,i)x_i \;+\; 2\sum_{p=1:n,p\ne i} A(i,p)x_p \notag \\
&=\; 2\sum_{p=1}^n A(i,p)x_p.
\end{align}

From (1.8), it follows that the $i$-th entry of the vector $2Ax$, where $1 \le i \le
n$, is given by
\begin{equation}
(2Ax)(i) \;=\; 2\sum_{p=1}^n A(i,p)x_p.
\end{equation}

From (10.53) and (10.54), we conclude that
\[
\pd{f}{x_i}(x) \;=\; (2Ax)(i), \quad \forall\, i = 1:n,
\]
and therefore, from (10.39), we obtain that
\begin{align*}
Df(x) \;&=\; \left(\pd{f}{x_1}(x) \quad \pd{f}{x_2}(x) \quad \ldots \quad \pd{f}{x_n}(x)\right) \\
&=\; \left((2Ax)(1) \quad (2Ax)(2) \quad \ldots \quad (2Ax)(n)\right) \\
&=\; 2(Ax)^t.
\end{align*}

To compute the Hessian of $f(x) = x^tAx$, let $j$ and $k$ fixed, with $1 \le j,k \le
n$. From (10.20), we find that
\begin{align*}
f(x) \;&=\; \sum_{1\le l,p\le n} A(l,p)x_lx_p \\
&=\; (A(j,k)+A(k,j))x_jx_k \;+\!\!\sum_{1\le l,p\le n;(l,p)\ne(j,k);(l,p)\ne(k,j)}\!\! A(j,k)x_lx_p \\
&=\; 2A(j,k)x_jx_k \;+\!\!\sum_{1\le l,p\le n;(l,p)\ne(j,k);(l,p)\ne(k,j)}\!\! A(j,k)x_lx_p,
\end{align*}
since $A$ is a symmetric matrix and therefore $A(j,k) = A(k,j)$. Then,
\[
\frac{\partial^2 f}{\partial x_k\partial x_j}(x) \;=\; 2A(j,k),
\]
and, from (10.41), it follows that
\[
D^2f(x) \;=\; \left(\frac{\partial^2 f}{\partial x_k\partial x_j}(x)\right)_{1\le j,k\le n}
\;=\; (2A(j,k))_{1\le j,k\le n} \;=\; 2A.
\]

(ii) Let $A$ be an $n \times n$ matrix. Recall from (10.28) that
\begin{equation}
q_A(x) \;=\; x^tAx \;=\; \frac12 x^t(A+A^t)x.
\end{equation}

Note that $A+A^t$ is a symmetric matrix, since
\[
(A+A^t)^t \;=\; A^t+(A^t)^t \;=\; A^t+A.
\]

Then, from (10.47) and (10.48), we find that
\begin{align}
D\left(x^t(A+A^t)x\right) &= 2\left((A+A^t)x\right)^t; \\
D^2\left(x^t(A+A^t)x\right) &= 2(A+A^t).
\end{align}

From (10.55), (10.56), and (10.57), we conclude that
\[
D(x^tAx) \;=\; \frac12 D\left(x^t(A+A^t)x\right) \;=\; \left((A+A^t)x\right)^t;
\]
\[
D^2(x^tAx) \;=\; \frac12 D^2\left(x^t(A+A^t)x\right) \;=\; A+A^t.
\qedhere
\]
\end{proof}

\textit{Vector Valued Functions}

Let $F : \R^n \to \R^m$ be a vector valued function given by
\[
F(x) \;=\; \begin{pmatrix} f_1(x) \\ f_2(x) \\ \vdots \\ f_m(x) \end{pmatrix},
\]
where $x = (x_1,x_2,\ldots,x_n)$.

If the functions $f_j(x)$, $j = 1:m$, are differentiable with respect to all
variables $x_i$, $i = 1:n$, then the gradient $DF(x)$ of the function $F(x)$ is the
following matrix of size $m \times n$:
\begin{equation}
DF(x) \;=\; \begin{pmatrix}
\dfrac{\partial f_1}{\partial x_1}(x) & \dfrac{\partial f_1}{\partial x_2}(x) & \ldots & \dfrac{\partial f_1}{\partial x_n}(x) \\[1.2ex]
\dfrac{\partial f_2}{\partial x_1}(x) & \dfrac{\partial f_2}{\partial x_2}(x) & \ldots & \dfrac{\partial f_2}{\partial x_n}(x) \\[1.2ex]
\vdots & \vdots & \ddots & \vdots \\[0.6ex]
\dfrac{\partial f_m}{\partial x_1}(x) & \dfrac{\partial f_m}{\partial x_2}(x) & \ldots & \dfrac{\partial f_m}{\partial x_n}(x)
\end{pmatrix}.
\end{equation}

If $F : \R^n \to \R^n$, then the gradient $DF(x)$ is an $n \times n$ square matrix.

Note that the $j$-th row of the gradient matrix $DF(x)$ is equal to the gradient
$Df_j(x)$ of the function $f_j(x)$, $j = 1:m$; cf. (10.39) and (10.58). Therefore,
\begin{equation}
DF(x) \;=\; \begin{pmatrix} Df_1(x) \\ Df_2(x) \\ \vdots \\ Df_m(x) \end{pmatrix}.
\end{equation}

\begin{lemma}
Let $M$ be an $m \times n$ matrix. The gradient of the vector valued function $Mx$,
taking $x \in \R^n$ into $Mx \in \R^m$ is
\begin{equation}
D(Mx) \;=\; M.
\end{equation}
\end{lemma}

\begin{proof}
Let $M = \mathrm{row}(r_j)_{j=1:m}$ be the row form of $M$. From (1.8), it follows
that the $Mx$ can be written as\footnote{Note that $r_jx$ is the result of the
multiplication of the row vector $r_j$ by the column vector $x$, and therefore a
real number, for every $j = 1:m$; see also (1.5).}
\begin{equation}
Mx \;=\; \begin{pmatrix} r_1x \\ r_2x \\ \vdots \\ r_mx \end{pmatrix}.
\end{equation}

From (10.59) and (10.61), it follows that
\begin{equation}
D(Mx) \;=\; \begin{pmatrix} D(r_1x) \\ D(r_2x) \\ \vdots \\ D(r_mx) \end{pmatrix}.
\end{equation}

From (10.43), we find that
\begin{equation}
D(r_jx) \;=\; r_j, \quad \forall\, j = 1:m.
\end{equation}

Then, from (10.62) and (10.63), we conclude that
\[
D(Mx) \;=\; \begin{pmatrix} r_1 \\ r_2 \\ \vdots \\ r_m \end{pmatrix}
\;=\; \mathrm{row}(r_j)_{j=1:m} \;=\; M.
\qedhere
\]
\end{proof}
